%% Cold Case: re-examining near-miss hits and flagged events of the
%% Breakthrough Listen nearby-star survey. AASTeX 6.3.1; pdflatex + bibtex.
%% Numbers: paper/paper_numbers.py. Figures: paper/make_figures.py.
%% Referee reports and responses: paper/REVIEW.md.
\documentclass[twocolumn]{aastex631}

\newcommand{\hzs}{Hz\,s$^{-1}$}
\newcommand{\sn}{S/N}

\shorttitle{Near-miss hits and flagged events of the BL 1327-star survey}
\shortauthors{Persechini}

\begin{document}

\title{Cold Case: Re-examining 771 Near-miss Hits and 34 Flagged Events from the Breakthrough Listen 1327-star Survey Using Public Data}

\author{Luiz Persechini}
\affiliation{Independent researcher}
%% \email{} % AUTHOR: fill in before submission

\begin{abstract}
Radio technosignature searches discard almost all of their detections with
automatic filters, and the physical nature of an individual discarded
detection is rarely examined. Using only public data, we re-measure 771
near-miss hits from the Breakthrough Listen search of 1327 nearby stars with
the Green Bank Telescope, chosen because they failed the survey's filter
narrowly, and the 34 events flagged ``Requires follow up'' in the survey's
public inspection table. Each case is measured again in the spectra of its
six-scan cadence against about 46 control windows and is closed by coded
rules that need direct evidence; the frequency band of a signal never closes
a case. We find no candidate. Of the 771 hits, 131 are shown to be local,
450 are single-sample bursts and not steady tones, 113 could not be
reproduced, 5 are carriers with sidebands, 2 have no public data, and 70
hits in 51 events are not explained. Of the 34 flagged events, 29 are
explained, 19 of them as measurable off the source; the survey's authors had
already rejected at least three of the 34 by cross-reference to other stars,
and our measurements agree. Hidden synthetic signals rode on every download:
the rules wrongly closed 4 of 918 recovered sky-only injections (0.4\%; 1.0\%
if injections masked by real signals are counted) and caught 72\% of weak
off-source leaks, but no injection drifts faster than 0.5\,\hzs, where three
quarters of the cases lie. We describe a frequency-permutation statistic for
the test of whether an event's line is seen toward other stars over days,
which needs no assumption about scan independence; it explains three flagged
events and shows that in crowded bands the test has little power. We also
document that the released hit files carry the drift-rate sign error noted
in the search software's documentation, characterize several interfering
clocks, and describe one unexplained event, toward HIP\,36357 at
2051.617\,MHz, whose frequency maps through the 240/221 transponder ratio to
within 0.2\,kHz of 2228.000\,MHz. A local or orbital transmitter is its
most likely origin. The code is public.
\end{abstract}

\keywords{Search for extraterrestrial intelligence (2127) --- Technosignatures (2128) --- Radio astronomy (1338) --- Astronomy data analysis (1858)}

\section{Introduction} \label{sec:intro}

Searches for narrowband radio technosignatures \citep{Tarter2001} produce
detections in very large numbers, and nearly all of them are radio-frequency
interference (RFI) from human technology. The Breakthrough Listen (BL)
program \citep{Worden2017,Isaacson2017,MacMahon2018} copes with this at the
Green Bank Telescope (GBT) by observing in cadences: three five-minute scans
of the target (ON) alternate with three scans of other positions (OFF). A
signal that is present in every ON scan and in no OFF scan is consistent
with a source fixed on the sky at the target's position
\citep{Enriquez2017,Price2020}. The \texttt{turboSETI} pipeline
\citep{turboseti} applies this as a filter on its list of hits, and what
survives is inspected by eye.

The filter is conservative by design. In the form used for the nearby-star
survey, a set of hits is discarded when any OFF hit lies within
$\pm600$\,Hz of the first ON hit, whether or not the two have the same
drift, and an event must be detected above threshold in every ON scan
\citep{Price2020}. \citet{Price2020} discuss the limitations of this
approach themselves, and its consequences for completeness have been
measured by injection and recovery: \citet{Margot2023} find that a process
imitating the BL pipeline recovers about 6\% of injected signals
\citep[see also][]{Margot2021}, and \citet{Pinchuk2019} discuss how much of
the band such blanking removes. The
same pipeline, with variations, underlies later BL searches
\citep[e.g.,][]{Gajjar2021,Choza2024}. Other work has re-searched the same
archive with different detection algorithms
\citep{Ma2023,Painter2025,JacobsonBell2025}.

This paper asks a narrower question about individual detections. For a
sample of hits that failed the filter narrowly, and for the events the
survey flagged for follow-up, what is each one? The filter is deterministic
and its decision can be reproduced from the hit tables; what is not recorded
is what the discarded detection physically was. We re-measure each case in
the public spectra and sort it into one of three classes: shown to be local
(a measurement that a source fixed on the sky cannot produce, such as the
same line with the telescope pointed elsewhere); shown not to be a steady
narrowband tone; or not reproduced. A case in none of them stays open, which
means only ``not explained by these checks''. A single observation can never
confirm a signal, so the product is the record of why each case was closed,
a short list of what was not, and measured error rates for the rules.

The project is the work of one person with a large-language-model coding
assistant on a desktop computer. It therefore leans on safeguards described
in Section~\ref{sec:method}: rules applied by code, hidden test signals
injected into every download, and cases with known answers. Before
submission the manuscript was also put through two adversarial reviews
against the code and the literature, which changed several conclusions; the
reports and responses are in the repository. The re-examination found no
candidate. Sections~\ref{sec:queue} and \ref{sec:followup} give the fate of
the 771 hits and the 34 flagged events; Section~\ref{sec:probe} describes
the statistic for the test over days; Sections~\ref{sec:emitters} and
\ref{sec:hip36357} describe the interfering clocks met along the way and
the one event we could not explain.

\section{Data} \label{sec:data}

\subsection{Hit tables and event tables}

\citet{Price2020} searched 1327 nearby stars between 1.10 and 3.45\,GHz and
released the products of the search \citep{PriceZenodo}. We use the
\texttt{turboSETI} hit files for the GBT L band (1.10--1.90\,GHz, with a
notch at 1.20--1.34\,GHz) and S band (1.80--2.80\,GHz, with a notch at
2.30--2.36\,GHz): 47.3 million unique hits in 13,286 scans, after removing a
duplicated copy of every scan. Each hit lists a frequency, a drift rate and
a signal-to-noise ratio (\sn). Only the strongest hit within $\pm600$\,Hz is
recorded \citep{Price2020}, so the absence of a catalogued hit is not the
absence of a line; this limits every use we make of the hit tables as
evidence of absence.

The same release holds the table of events that the survey inspected by eye:
6094 at the GBT (4522 at L band and 1572 at S band, the event groups of
Table~4 of \citealt{Price2020}). The table has a label field. It is
``None'' for 5363 events (88\%); 583 are labelled ``Hits in off'', 64
``Plotting issue'', 50 ``Interesting but not ET'', and 34 ``Requires follow
up'' (32 at L band and 2 at S band). We treat the label as a working
annotation and not as the survey's final judgment: \citet{Price2020} report
that every event group was inspected and that none survived, and their
Appendix~A discusses and rejects, by cross-reference to similar signals
toward other stars, at least three of the 34 (the events of their Figures
7a--7c, toward HIP\,54677 at 1435.943 and 1439.068\,MHz and toward
HIP\,100064 at 1595.878\,MHz). \citet{Painter2025} later reached the same
verdict on the HIP\,54677 and HIP\,100064 signals.

\subsection{A note on the sign of the drift rate} \label{sec:sign}

In the released hit files the listed drift rate has the opposite sign to
the motion of the line in the data: a line listed with a positive drift
moves down in frequency from one scan to the next. The tables show this
without raw data. We paired hits of the same line in two ON scans of one
cadence (the first and third, or the third and fifth) and asked which sign
of the listed drift predicts the second hit's frequency. Pairs were kept
when the first hit had \sn\ of 20 or more, a drift between 0.05 and
1\,\hzs, and a unique match within 1\,kHz. Of 1,129,865 pairs, the frequency
of the second hit matched the prediction of the reversed sign in 1,024,573
and that of the listed sign in 12,246; the rest matched neither. That is
98.8\% of the decided pairs.

The cause is documented. The \texttt{turboSETI}
documentation\footnote{\url{https://github.com/UCBerkeleySETI/turbo_seti/blob/master/docs/find_event.rst}}
notes that hit files produced before version 0.8.2 of the software and
version 1.1.7 of \texttt{blimpy} (about mid-2019) recorded drift rates with
the incorrect sign, to be flipped when events are built. The survey did so:
in the released event table the drift rates have the opposite sign to the
hit files in every one of 2809 entries we compared, so the survey's results
are unaffected. The hit files themselves are not flagged, and the record
does not state the software version. Anyone who extrapolates a hit to
another scan from these files alone, as our pipeline first did
(Section~\ref{sec:errors}), follows every drifting line the wrong way.
\citet{Margot2023} report a similar sign reversal in the drift rates of
\citet{Ma2023}, which derive from the same archive.

\subsection{Spectra}

The raw data are the public high-spectral-resolution filterbank products
\citep{Lebofsky2019}: about 2.79\,Hz per channel and about 18\,s per time
sample, 16 samples in a five-minute scan. One scan is about 15\,GB. The
files are HDF5 with one storage chunk per coarse channel (2.9\,MHz) and time
sample, and the archive honours HTTP range requests, so the pipeline reads
only the chunks around a hit: about 60\,MB per scan. No complete file was
downloaded.

\section{Method} \label{sec:method}

\subsection{Selection of cases} \label{sec:selection}

From the hit tables we built tracks: a hit in an ON scan followed along its
drift into the other five scans of its cadence. We selected tracks with a
drift-consistent hit in only one or two of the three ON scans and no OFF hit
on the track. Cuts kept tracks inside the receiver band and outside its
notch, with $|{\rm drift}|<3.9$\,\hzs, with no hit toward any other target
of the survey within about 1\,kHz, from scans that contribute at most 20
such tracks, and with \sn\ of 25 or more. That gave 781 tracks. Nineteen
were used in a pilot in which the first thresholds were set and are not
counted further; 762 entered the queue. When the selection was repeated with
the drift sign corrected (Section~\ref{sec:sign}), 740 of the 781 were
selected either way, 41 dropped out (they were kept) and 9 were new and
were added, for 771 cases toward 296 stars.

In 697 of the 771 no OFF hit lies within 1.2\,kHz of the track at all; in 74
one does, but off the track. Our 1.2\,kHz is the window of the current
software and is twice the $\pm600$\,Hz the survey used, so the 697 are clean
by a stricter standard than the survey's and failed only its requirement of
a hit in every ON scan.

The sample is not a random sample of the discarded hits, and two of the
cuts shape the results. The median drift is 1.7\,\hzs; only 192 cases drift
by 0.5\,\hzs\ or less. And the cut on hits toward other targets removes by
construction the lines that the tables show toward other stars, so a later
finding that the hit tables hold nothing near a queue case toward other
stars is largely guaranteed and is not evidence.

The 34 flagged events were entered separately, each at its strongest hit,
and run through the same code unchanged. None of them is among the 771.

\subsection{Forced measurement against control windows}

All \sn\ values of this work are on one scale, which is not that of the hit
tables. A cutout is divided by its bandpass (a running median along
frequency), summed along a trial drift, and expressed as the deviation from
the median of the summed spectrum in units of 1.4826 times its median
absolute deviation. On this scale the HIP\,36357 line of
Section~\ref{sec:hip36357} has \sn\ 75 and 134 where the hit table lists 17
and 30.

For each case the pipeline first finds the catalogued hit in the data (the
seed): the strongest line within 50\,Hz of the catalogue position, plus the
drift over one time sample, at either sign of the catalogue drift; every ON
scan is searched and the strongest kept. It then measures every one of the
six scans at the place predicted from the seed, whether or not the hit
table lists a hit there. The measurement is the largest \sn\ over start
frequencies within a tolerance of the predicted one (20\,Hz, growing with
the time from the seed scan and with the drift) and over drifts within
$\pm(0.05+0.25|d|)$\,\hzs\ of the seed's drift $d$. The same maximum is
taken in 46 control windows of the same width at other frequencies of the
cutout. A track is present in a scan when it exceeds the strongest control
window by 1.5. The three OFF scans are also stacked. Taking the strongest
control as the reference is conservative in crowded spectra, at the cost of
missing lines there; that cost recurs throughout this paper.

\subsection{Rules that close a case} \label{sec:rules}

A case is closed by code when one of the kinds of evidence in
Table~\ref{tab:rules} is found. The rules fall into three classes, and the
distinction matters for what a closure means. Rules D1, D3, D5, D7 and D9
show a line with the telescope pointed elsewhere, and D2 and D6 show the
arithmetic of a human clock: these are evidence of a local origin. Rule D4
shows that the catalogued hit is not a steady tone: its power lies in one
18\,s sample. A source on the sky could in principle pulse, so D4 is not
evidence of locality; it establishes that a de-Doppler search over a
five-minute scan reported a short burst as a drifting tone. A case with no
line at the catalogue position is ``not reproduced''.

Where a rule depends on a coincidence, such as a line near the predicted
place in an OFF scan, it closes the case only when the chance of an
unrelated line doing the same is below 5\%, with that chance estimated from
the same data and never taken below 1\% per scan. The allocation or
occupancy of the band is never used to close a case.

\begin{deluxetable*}{llp{11.5cm}}
\tabletypesize{\footnotesize}
\tablecaption{Evidence that closes a case \label{tab:rules}}
\tablehead{\colhead{Rule} & \colhead{Evidence} & \colhead{Meaning}}
\startdata
D1 & Seen off the source & The line is measurable in the OFF scans: on its predicted track in two scans or in their stack; as a line of the same drift and \sn\ of 8 or more within 300\,Hz to 1\,kHz; or by following a line whose drift changes \\
D2 & A comb at a round rate & Companion lines spaced by a round number of hertz \\
D3 & Other stars, hit tables & The same frequency and drift toward two or more other stars in the hit tables \\
D4 & A burst, not a tone & Half or more of the track's power is in one 18\,s sample, or nothing is left of the track without its strongest sample \\
D5 & Another day & The line is back at the same topocentric frequency when the star was observed again \\
D6 & Harmonics of one clock & Lines of one scan at whole multiples of one clock rate, all the same fraction off nominal \\
D7 & The same half hour & The line is within 300\,Hz in scans of other stars, at least 1 degree away, within 30 minutes \\
D8 & An explained companion & A line that rises and falls with the event in the same scans is itself closed by D1, D4, D5 or D7; or the event and two such lines are equally spaced, as a carrier and its sidebands \\
D9 & The days around & The line is at the event's topocentric frequency toward other stars over days, more consistently than at other frequencies of the same spectra (Section~\ref{sec:probe}) \\
\enddata
\tablecomments{D1--D4 are applied automatically and are the rules exercised by the hidden test signals; D5--D9 run afterwards on what is still open. D3 can hardly fire on the queue, because of the selection (Section~\ref{sec:selection}).}
\end{deluxetable*}

Rule D5 has a counterpart that is never decided by code: a line that
reappears at the frequency shifted by the Earth's motion, where a
transmitter fixed in the barycentric frame would have moved, is flagged for
a person. Not seeing a line again changes nothing.

\subsection{Hidden test signals} \label{sec:controls}

Injection and recovery is the standard way to measure a search's efficiency
\citep{Margot2021,Brzycki2022}. Here 476 of the 485 download groups
received three synthetic narrowband signals each, injected into the cached
data before measurement and passed through the same code as the real cases.
Two are present in ON scans only, in all three scans or in one or two of
them, and should stay open. The third is also present in the OFF scans at
12\% to 25\% of its ON strength (a weak leak) and should be closed. The
nominal injected \sn\ is 12 to 60 for the sky-only signals and 25 to 80 for
the leaks; drifts are 0.01 to 0.5\,\hzs\ with a small curvature, and the
signals scintillate between scans and flicker within them. They sit at
random frequencies at least 20\,kHz from the real cases. An injection is set
aside as ``masked'' when what was measured at its position is a real signal
more than three times stronger or of a different drift; that category was
defined after the first failures were inspected.

The test signals have limits that matter for every rate quoted below. They
are injected with the pipeline's own conventions, so they cannot reveal a
misreading of the catalogue (Section~\ref{sec:errors}). They exercise rules
D1 to D4 only. And none drifts faster than 0.5\,\hzs, while three quarters of
the queue and 32 of the 51 open events do.

\subsection{Other stars in the days around} \label{sec:probe}

A local emitter does not know where the telescope points, so the same line
toward a different target is evidence of a local origin. The idea is old.
The Allen Telescope Array searches used anticoincidence between
simultaneous beams \citep{Harp2016}; direction-of-origin filters compare
detections across pointings \citep{PinchukMargot2022,Margot2021};
\citet{Price2020} rejected their best events by cross-reference to similar
signals toward other stars, as did \citet{Painter2025}; and an archival
search at the same frequency is a step of the verification framework of
\citet{Sheikh2021}. What we add is a forced measurement in the spectra of
other targets' scans over days, with a chance estimate that does not assume
those scans are independent.

For an event at topocentric frequency $f_e$ and drift $d_e$, the pipeline
lists the scans of other stars in the same receiver band within seven days,
nearest in time first, excluding the event's own cadence, and reads a cutout
of at least $\pm355$\,kHz around $f_e$ from each until a 3\,GB cap (30 to
108 scans here). The Earth-shifted frequency is not measured, because it
has no meaning toward a different star.

Our first statistic for these data was wrong, and the way it failed is
instructive. It asked in how many scans the event's window exceeded its
strongest control window, and took the chance of that from null frequencies
chosen where the hit tables list nothing. The review found three faults:
nulls chosen to be clean bias the chance low, since the event's frequency
is by construction not clean; where the control windows hold strong lines
neither the event nor a null can ever match, so ``no null matched'' carries
no information; and the nulls' control windows overlapped the event's
window and each other. It produced probabilities as small as $10^{-135}$.

The statistic we use instead compares the event's frequency with every
other frequency of the same cutouts, taken as they are. Each cutout is laid
out in adjacent windows of half-width $W$, one of them centred on $f_e$. In
scan $s$ the strongest line of drift within $\pm(0.05+0.25|d_e|)$\,\hzs\ of
$d_e$ is found in each window, and the event's window is ranked among the
$N$ others: $p_s$ is the fraction at least as strong, counting itself. The
statistic for a window is $T=-\sum_s\ln p_s$, computed in the same way for
every window, and
\begin{equation}
P=\frac{1+\#\{j: T_j\ge T_e\}}{1+N}
\end{equation}
is the fraction of windows that show a line as consistently as the event's
does. A persistent unrelated line earns a high $T$ at its own window, as it
should, so scans need not be independent; nothing is conditioned on the hit
tables; and no control window or threshold enters. $P$ cannot be smaller
than $1/(1+N)$. We use $W=50$, 500 and 2000\,Hz, the last for lines that
wander between nights, and multiply the smallest $P$ by three. Rule D9
closes an event when this corrected $P$ is below 0.05 and its window is in
the top 1\% of its scan in at least two scans.

The statistic was designed after the data were seen, in response to the
review, and the three widths were informed by the offsets observed, so it is
not a pre-registered test. Its null hypothesis is that the event's
frequency is unrelated to the lines seen toward other stars; under it, $f_e$
is exchangeable with the other windows of the cutout, which holds for a sky
signal whatever the local RFI. It was run on nine events, so a further
factor of nine applies to any one of them taken as a discovery.

\subsection{Use of language models} \label{sec:ai}

Two language-model readers (Anthropic's Claude Sonnet) wrote up and rated
each open event. They rank; they cannot close a case. That boundary was set
after a test: asked what an injected sky-only signal was ``most likely'' to
be, the readers called 12 of 21 interference, arguing from the band and its
surroundings; under the instruction to close only on direct evidence they
kept 20 of 21. The readers also failed to connect lines that belonged to one
device (Section~\ref{sec:emitters}). The pipeline, its tests and a first
draft of this text were written with a language-model coding assistant, and
the two reviews mentioned above were made by further instances of such
models (Section~\ref{sec:ack}).

\section{Validation} \label{sec:validation}

\subsection{Known answers}

The Voyager~1 carrier at 8419.297\,MHz, in the public GBT observation of
2016 September 19 that accompanies the \texttt{turboSETI} software, is left
open, with a measured drift of $-0.37$\,\hzs\ against $-0.374$\,\hzs\
predicted from the Earth's motion alone. The signal of interest blc1 toward
Proxima Centauri \citep{Smith2021}, which \citet{Sheikh2021} identified as
an intermodulation product of local oscillators, is closed: five other lines
in the same data are spaced as multiples of a 2\,MHz clock running 2.5\,ppm
fast.

Of the eight signals of interest of \citet{Ma2023}, four are closed as
measurable in the OFF scans (MLc3, MLc5, MLc6 and MLc7) and four are left
open (MLc1, MLc2, MLc4 and MLc8). \citet{Margot2023} had already examined
these: they detect MLc3, MLc4, MLc7 and MLc8 in the OFF scans and attribute
MLc5 to interference by its similarity to MLc4. We agree on MLc3, MLc5 and
MLc7, and our off-source test misses two (MLc4 and MLc8) that theirs finds,
which is a measured example of its limited sensitivity. MLc6 lies 2.4\,kHz
from the HIP\,54677 line at 1435.943\,MHz.

\subsection{Injections}

Table~\ref{tab:controls} gives the outcome of the 1428 hidden test signals
of the queue, and Figure~\ref{fig:controls} shows it against injected
strength and drift. Of the 952 sky-only injections, 914 were left open, 24
were not recovered, 10 were masked by a stronger real signal, and 4 were
recovered and closed by a rule: 4 of 918, or 0.4\% (95\% interval 0.2\% to
1.1\%). Three of the four were closed as seen off the source, in crowded
spectra where real interference sat near the predicted track in an OFF
scan. One of the unrecovered signals was also closed, as a burst, under a
real burst ten times its strength. In real data one cannot know that a
signal is masked: counting the 5 masked signals that were closed, the rate
is 9 of 928, or 1.0\% (0.5\% to 1.8\%). Counted unconditionally, 33 of 952
sky-only signals (3.5\%) did not stay open, most of them because they were
too weak to recover.

Of the 469 unmasked leaks, 340 were caught (72\%; 68\% to 76\%). That
figure describes the chosen mix of injections more than the rules. The
catch rate rises with the nominal off-source \sn\ from 8\% below 4, through
34\% (4 to 6) and 64\% (6 to 8), to 92\% (8 to 12) and 99\% above 12; and
it falls with drift, from 95\% below 0.1\,\hzs\ to 47\% at 0.4 to
0.5\,\hzs\ (Figure~\ref{fig:controls}). Recovery of sky-only signals also
falls at the fastest drifts tested, and 22 of the 24 unrecovered signals
drift faster than 0.35\,\hzs. Extrapolated to the 1 to 4\,\hzs\ of most of
the queue, the off-source test probably has little sensitivity, so ``not
seen off the source'' carries little weight for the fast-drifting open
events, and the false-closure rate of the rules there is unknown. A missed
leak stays on the open list, so this error inflates the open list and does
not remove signals.

These rates are in-sample. The count of false closures was 13 before the
controls were examined one by one; eight had been closed by rule D3 on a
single weak catalogued hit toward one other star, which is why that rule
now needs two stars, and the masked category removed others. The only fresh
sample is the set of test signals carried by the flagged events' downloads:
none of 59 recovered sky-only signals was closed (95\% upper limit about
6\%), and 14 of 30 leaks were caught, the missed ones being weak off the
source or in crowded spectra.

\begin{deluxetable}{lrrr}
\tablecaption{Outcome of the hidden test signals \label{tab:controls}}
\tablehead{\colhead{Set} & \colhead{Sky-only} & \colhead{Closed} & \colhead{Leaks} \\
\colhead{} & \colhead{left open} & \colhead{by a rule} & \colhead{caught}}
\startdata
Pilot & 20 of 21 & 0 & 9 of 10 \\
Holdout, random frequencies & 18 of 20 & 0 & 4 of 10 \\
Queue, running & 914 of 952 & 4\tablenotemark{a} & 340 of 476 \\
Flagged events & 59 of 60 & 0 & 14 of 30 \\
\enddata
\tablenotetext{a}{Of 918 recovered and unmasked. A further 24 were not recovered (one of them closed as a burst) and 10 were masked (5 of them closed). Seven of the 476 leaks were masked.}
\end{deluxetable}

\begin{figure*}
\plotone{figures/controls.pdf}
\caption{Fraction of the hidden test signals of the queue decided as
intended, against nominal injected \sn\ (left) and against drift (right),
in bins of equal count: sky-only signals left open (blue; unrecovered
signals count as failures) and weak leaks closed (red). Masked injections
are left out. Error bars are 68\% Wilson intervals. No injection drifts
faster than 0.5\,\hzs. \label{fig:controls}}
\end{figure*}

\subsection{Two errors of our own} \label{sec:errors}

Two errors were found by hand, not by the controls, and every stored case
was measured again from cached data after each.

Until the drift sign of Section~\ref{sec:sign} was noticed, the pipeline
followed every drifting case in the wrong direction. A wrong-way track
crosses a steady line for an instant, so steady lines looked like bursts.
After the correction the seed moved to another track in 610 of 761 cases
and 125 changed status, in both directions (66 that had been closed as
bursts were reopened or closed for another reason; 13 open cases were
closed).

The second error was a seed search that looked too far from the catalogue
position for fast drifters. At 484 places with nothing catalogued it
returned a seed 43\% of the time. Restricted to 50\,Hz plus the drift over
one time sample, it returns one 14\% of the time, mostly from bursts, and
recovers injected fast tones slightly more often.

\section{Results} \label{sec:results}

\subsection{The 771 near-miss hits} \label{sec:queue}

Table~\ref{tab:queue} gives the outcome by class. None is a candidate.

\emph{Shown to be local: 131.} In 102 cases the line is measurable in the
OFF scans once the predicted track is measured there directly. Twenty-three
are lines of clocks (Section~\ref{sec:emitters}). Five were found toward
another star within half an hour, and one moves with a companion that is
seen off the source.

\emph{Not a steady tone: 450.} In 243 cases the track stood above its
controls over the whole scan and does not without its strongest 18\,s
sample. In 202 the rule fired only on its other condition, half or more of
the power in one sample; in 84 of those the track still stands above its
controls without that sample, so these are steady or slowly varying lines
with a burst on top, or strongly varying lines, and they are the weakest of
the closures (73 have a share between 50\% and 60\%). Five move with a
companion that is itself a burst. Most of these cases drift fast: 85\% of
the burst closures have $|{\rm drift}|>0.5$\,\hzs, outside the range of the
test signals.

\emph{Not reproduced: 113.} In 33 cases no line was found at the catalogue
position. In 80 more no line was found and the cutout holds a strong
single-sample burst, which is the likely origin of the catalogued hit. We
have not measured how often the seed search recovers an ordinary catalogued
hit of this strength at these drifts, so a failure to reproduce is not
evidence that the hit was spurious.

\emph{Carriers with sidebands: 5.} Five were closed by the clause of rule
D8 that takes three equally spaced lines rising and falling together as a
carrier and its sidebands. That is evidence of modulation and not of a
local origin (the Voyager~1 signal has sidebands), and we list these
separately.

\emph{Open: 70 hits in 51 events.} By the criterion of
Section~\ref{sec:method}, 30 events are present in one ON scan, 2 in two,
none in all three, and 19 in none: in those 19 the line was found by the
seed search but does not exceed its strongest control window by 1.5 in any
scan. Twenty-seven of the 51 drift faster than 1\,\hzs.

Three further checks were run on the open events. For 18 events the star
was observed again in the same band, on 27 other days in all, and the line
was absent every time at both the topocentric and the Earth-shifted
frequency. A comparison of each line of sight with the orbits of catalogued
satellites flagged about as many close passes as chance predicts and closes
nothing. The half-hour test (D7) closed five hits, and these closures are
weak: with the 1\% floor on the per-scan chance, one matching scan closes a
case exactly when five or fewer scans were available, so the quoted chance
of 4\% to 5\% is arithmetic and not a measurement; the single matches lie
138 to 304\,Hz from the event, two of them at the edge of the 300\,Hz
window, where the matches of Section~\ref{sec:days} lie within tens of
hertz; and with some sixty events tested, about three such matches are
expected by chance. They should be read as flags.

\begin{deluxetable}{lr}
\tablecaption{Outcome for the 771 near-miss hits \label{tab:queue}}
\tablehead{\colhead{Outcome} & \colhead{Hits}}
\startdata
\multicolumn{2}{l}{\emph{Shown to be local}} \\
\quad Seen off the source (D1) & 102 \\
\quad Lines of human clocks (D2, D6) & 23 \\
\quad Toward another star within half an hour (D7) & 5 \\
\quad Companion seen off the source (D8) & 1 \\
\multicolumn{2}{l}{\emph{Not a steady narrowband tone}} \\
\quad Nothing left without the strongest sample (D4) & 243 \\
\quad Half or more of the power in one sample (D4) & 202 \\
\quad Companion is a burst (D8) & 5 \\
\multicolumn{2}{l}{\emph{Not reproduced}} \\
\quad No line; a strong burst in the cutout & 80 \\
\quad No line at the catalogue position & 33 \\
\multicolumn{2}{l}{\emph{Other}} \\
\quad Carrier with sidebands (D8) & 5 \\
\quad No public data & 2 \\
\quad Not explained by these checks & 70 \\
\enddata
\end{deluxetable}

\subsection{The 34 flagged events} \label{sec:followup}

Table~\ref{tab:followup} gives the outcome for the events labelled
``Requires follow up''. Nineteen are measurable in the OFF scans when the
predicted track is measured there directly; two of those rest on the
stacked OFF scans alone, at a margin of 1.2 over the strongest control where
1.0 is required. Three are bursts. Two, both near
1420.4\,MHz at drifts of about $-1.2$\,\hzs, have no line at the catalogue
position; they lie on the neutral-hydrogen line, where the emission differs
between ON and OFF positions, and at a drift where our recovery is
untested. One event (HIP\,109427 at 1554.807\,MHz) has its line toward
other stars in all five scans of the same half hour, at the same frequency.
One (1691.081\,MHz toward HIP\,17420) has three lines 45 to 172\,kHz below
it that switch on and off with it (correlation 0.97 to 0.99 between time
samples), each at \sn\ above 400 and each seen off the source.

\begin{deluxetable}{lr}
\tablecaption{Outcome for the 34 flagged events \label{tab:followup}}
\tablehead{\colhead{Outcome} & \colhead{Events}}
\startdata
Seen off the source (D1) & 19 \\
A burst, not a tone (D4) & 3 \\
No line at the catalogue position & 2 \\
Toward other stars within half an hour (D7) & 1 \\
Moves with an explained companion (D8) & 1 \\
Toward other stars in the days around (D9) & 3 \\
Not explained by these checks & 5 \\
\enddata
\end{deluxetable}

\subsection{Other stars in the days around} \label{sec:days}

Table~\ref{tab:probe} and Figure~\ref{fig:probe} give the statistic of
Section~\ref{sec:probe} for the nine events probed: the eight flagged
events that the automatic rules had left open, and the HIP\,36357 event.

Three flagged events meet the criterion. For HIP\,100064 at 1595.877\,MHz
the line is in the top 1\% of its scan in 43 of 54 scans of other stars
over the following 27 hours, moving through the 2\,kHz window within each
observing session, toward 28 stars; $T$ is nearly twice the largest value among the 174
other windows, and $P$ is at the floor that number sets. For HIP\,54677 at
1435.943\,MHz the line is within 40\,Hz of the event's frequency toward
seven other stars, mostly 10 to 11 hours earlier; none of 7150 other windows
reaches its $T$. These are the two events that \citet{Price2020} rejected on the
same grounds from the hit tables; the measurement in the spectra confirms
their verdict and adds how long and toward how many stars the lines
persist. The third, HIP\,100064 at 1560.945\,MHz, has its line within
51\,Hz in seven scans of four other stars 27 to 28.5 hours later, at \sn\
up to 21, in spectra crowded enough that seven of 7150 other windows do as
well.

For the other five flagged events the criterion is not met, and for three
of them that is a statement about the method and not about the sky. Toward
HIP\,109427 at 1559.671\,MHz a line of the event's drift is within 17\,Hz
in eleven scans of other stars in the hour and a half before, at \sn\ of 67
to 448; but 2.7\% of the other 100\,Hz windows of those spectra hold a line
as consistent, and the corrected $P$ is 0.08. Toward the same star at
1560.003\,MHz there is one at 25 to 28\,Hz in two scans ($P=0.16$). Toward
HIP\,29432 at 1556.346\,MHz there is one in ten consecutive scans, sliding
from 78 to 117\,Hz at \sn\ of 35 to 92, in spectra where a third of the
windows of that width hold something comparable. In these bands a line near
any chosen frequency is common, so its presence cannot be told from chance
and the events stay open. For HIP\,42173 and HIP\,87108 nothing stands
out. Only for HIP\,36357 is absence informative: its spectra are clean
(strongest control window below \sn\ 6 in every scan), and its line is
not there.

Four of the five events left open lie between 1538 and 1561\,MHz and drift
by 0.02\,\hzs\ or less.

\begin{deluxetable*}{llrrrrl}
\tabletypesize{\footnotesize}
\tablecaption{The line at the event's topocentric frequency in scans of other stars within seven days \label{tab:probe}}
\tablehead{\colhead{Event} & \colhead{Frequency} & \colhead{Scans} & \colhead{$W$} & \colhead{Scans in} & \colhead{Corrected} & \colhead{Outcome} \\
\colhead{} & \colhead{(MHz)} & \colhead{} & \colhead{(Hz)} & \colhead{top 1\%} & \colhead{$P$} & \colhead{}}
\startdata
HIP\,100064 & 1595.877 & 54 & 2000 & 43 & 0.017\tablenotemark{a} & Local \\
HIP\,54677 & 1435.943 & 30 & 50 & 11 & 0.0004\tablenotemark{a} & Local \\
HIP\,100064 & 1560.945 & 108 & 50 & 10 & 0.003 & Local \\
HIP\,109427 & 1559.671 & 30 & 50 & 7 & 0.08 & Open; a line within 17\,Hz in 11 scans \\
HIP\,109427 & 1560.003 & 30 & 50 & 2 & 0.16 & Open; a line at 25--28\,Hz in 2 scans \\
HIP\,29432 & 1556.346 & 47 & 500 & 0 & 0.96 & Open; a line at 78--117\,Hz in 10 scans \\
HIP\,42173 & 1538.712 & 54 & 500 & 0 & 1 & Open \\
HIP\,87108 & 2199.097 & 30 & 2000 & 0 & 1 & Open \\
HIP\,36357 & 2051.617 & 54 & 500 & 1 & 0.49 & Open; clean spectra (Section~\ref{sec:hip36357}) \\
\enddata
\tablenotetext{a}{At the floor set by the number of windows: no other window reaches the event's $T$.}
\tablecomments{$W$ is the window half-width that gives the smallest $P$; the corrected $P$ is three times that value. The first eight rows are flagged events; the last is from the queue.}
\end{deluxetable*}

\begin{figure*}
\plotone{figures/probe.pdf}
\caption{Left: the event toward HIP\,100064 at 1595.877\,MHz in 54 scans of
35 other stars after the event: offset of the strongest line of the event's
drift within 2\,kHz. Filled points are scans in which that window is in the
top 1\% of the 2\,kHz windows of its spectrum. Right: the statistic $T$ of
Section~\ref{sec:probe} for the nine events probed (red), against the other
windows of the same spectra: the bar spans their median to their 99th
percentile, and the tick marks their maximum. \label{fig:probe}}
\end{figure*}

The test signals carried by the flagged events' downloads were decided
less well than in the queue (Table~\ref{tab:controls}), so ``open'' is
weaker evidence here than there.

\subsection{Where the flagged events were observed} \label{sec:belt}

The flagged events are not spread over the survey. Of the 22 L-band stars
with a flagged event, 19 were pointed within 13 degrees of the
geostationary belt as seen from Green Bank (median 7.3 degrees), against
33\% of the 890 L-band stars of the survey (median 26 degrees;
Figure~\ref{fig:belt}). They are also clustered in time: 18 of the 22 were
observed between MJD 57510 and 57600 (2016 May to July), when 34\% of the
L-band stars were.

We report this as a pattern and attach no probability to it. The 13 degree
threshold was chosen after looking: a binomial comparison with the survey
rate gives values from 0.06 at 5 degrees to $10^{-7}$ at 15 degrees. From
Green Bank a pointing within 13 degrees of the belt is a pointing at
declination $-19$ to $+7$ degrees, so angle from the belt and declination
are the same variable. And the stars share observing sessions, so they are
not independent. What the data support is that the events flagged for
follow-up come mostly from low-declination targets observed in a few months
of 2016. An emitter on or near the geostationary belt would produce that,
and so would a ground source toward the south or one active in that season.
\citet{Price2020} already associate the hit density at 1.525--1.560\,GHz
with satellite downlinks. The belt does not single out the open events: the
four between 1538 and 1561\,MHz lie 6 to 13 degrees from it, no closer than
the closed events of the same band. That range straddles 1559\,MHz, the
boundary between the mobile-satellite downlink allocation and the
radionavigation-satellite allocation, and a satellite transmitter is a
natural hypothesis for nearly undrifting lines there; we did not test it
against the positions of individual satellites, and by our rule a band does
not close a case.

\begin{figure}
\plotone{figures/belt.pdf}
\caption{Angle between the pointing and the geostationary belt as seen from
Green Bank, for all L-band stars of the survey (grey) and for the L-band
stars with an event flagged for follow-up (red). Ticks mark the stars of
the flagged events still open. \label{fig:belt}}
\end{figure}

\subsection{Interfering clocks} \label{sec:emitters}

Several groups of hits trace to one device each (Table~\ref{tab:emitters}).
We can characterize these devices but not place them: they may be at the
site, on a satellite, or in the instrument itself.

Eight hits in one scan toward HIP\,17420 are lines spaced by exactly
131.072\,kHz, which is $2^{17}$\,Hz; in that scan, 24 of the 29 catalogued
hits within 60 teeth of the comb sit on a tooth, and none do when the grid
is shifted. Eleven events in five scans toward three stars, between 2016
February 29 and March 3, sit at 59, 62, 69 and 77 times 24\,MHz, every line
of a scan the same fraction below nominal to 0.1\,ppm or better: a 24\,MHz
clock running 70 to 103\,ppm slow, or a submultiple of it. A twelfth event,
toward a fourth star (HIP\,87108) on March 3, sits at 59 times 24\,MHz with
companions at 58 and 62 times, all 108\,ppm slow. The rule tries every
whole number of megahertz from 4 to 50 with an allowance of 200\,ppm, so a
single line is near some multiple by chance; it therefore needs two lines
of one scan that agree to 0.3\,ppm, and it was checked on predictions it had
not used: at three other multiples per scan the spectra hold a line at the
predicted frequency in 13 of 15 measurements, and the hit tables hold a hit
at 48 of 144 other predicted multiples against 3 of 2880 shifted control
positions. A weak line at 2258.417\,MHz toward HIP\,15919 is one tooth of a
comb of far stronger lines spaced by 2.048031\,MHz, the rate of the E1
telecommunication clock and of common audio converters, running 15\,ppm
fast.

The language-model readers had read eight of the 24\,MHz lines one at a
time and left all eight open. Connecting lines that are tens of megahertz
apart, in different downloads, is a job for a rule that looks across cases.

\begin{deluxetable*}{lll}
\tabletypesize{\footnotesize}
\tablecaption{Interfering devices characterized \label{tab:emitters}}
\tablehead{\colhead{Signature} & \colhead{Where and when} & \colhead{Evidence}}
\startdata
Comb at 131.072\,kHz ($2^{17}$\,Hz) & L band, toward HIP\,17420 & 24 of 29 catalogued hits on a tooth \\
24\,MHz clock, 70--108\,ppm slow & L band, 2016 Feb 29--Mar 3, four stars & Harmonics 58 to 77; 13 of 15 predicted lines found \\
Comb at 2.048031\,MHz (15\,ppm fast) & S band, toward HIP\,15919 & Lines at 1, 2, 3 teeth, \sn\ up to 55,742, on and off source \\
Drifting line near 1595.88\,MHz & L band, 27 hours, 28 stars & Table~\ref{tab:probe}; also \citet{Price2020} \\
Pulsed lines 6\,kHz apart near 2051.62\,MHz & S band, 2018 Apr 15 & More than twenty lines in one OFF scan (Section~\ref{sec:hip36357}) \\
\enddata
\end{deluxetable*}

\subsection{An event not yet explained: HIP 36357 at 2051.617 MHz} \label{sec:hip36357}

One event of the queue stands apart (Figure~\ref{fig:hip36357}). On 2017
February 23, between 01:27 and 02:01 UT, a narrow line toward HIP\,36357 at
2051.616734\,MHz appears in the first two ON scans, at \sn\ of 75 and 134,
and in none of the OFF scans (\sn\ 2.5 to 3.1, with controls up to 4.8). It
is absent from the third ON scan. The two detections lie on one straight
track of $-0.124$\,\hzs\ to within 0.5\,Hz over 12 minutes. The hit table
lists it in both of those scans, at \sn\ 17 and 30 and frequencies 89\,Hz
apart, and not in the third, so a filter that requires a hit in every ON
scan does not retain it.

The line switches on and off by itself. In the first ON scan it is on from
the first sample, varying tenfold, and off in the last two samples. In the
second ON scan it is off for the first 110\,s and then on to the end,
rising. The first OFF scan falls entirely inside a gap of the line, and the
second began 45\,s after it was last seen, when it had fallen from \sn\ 134
to below 4.8. Its absence off the source therefore does not show that it
comes from the sky, and we do not count it as evidence either way.

Checks that found nothing, with their limits. The star was observed again
three days later, in a single five-minute scan, and the line was absent at
the topocentric frequency and at the Earth-shifted one (8.75\,kHz lower).
It was absent toward the one other star observed within 30 minutes; that
star, HIP\,36366, is 12.6 arcminutes away, about two beam widths, so for an
emitter in a far sidelobe it is the same pointing at another time and not
an independent direction. The nearest catalogued satellite at the strongest
moment was 4.0 degrees from the line of sight, which is ordinary for that
night. The line was absent in 54 scans of 36 other stars from 31 hours
before to 5.5 hours after, in clean spectra (Table~\ref{tab:probe}). The
hit tables list nothing within 2\,kHz toward other stars within 30 days,
which the selection largely guarantees. A transmitter at rest in the
barycentric frame would drift at $-0.179$\,\hzs\ in that direction; the
line drifts at $-0.124$, which neither supports nor excludes a sky origin,
since the source may have an acceleration or an oscillator drift of its
own.

Two leads in the hit tables point to human transmitters. Three catalogued
lines of the same drift 128\,MHz below the event turned out, in the
spectra, to belong to a forest of several hundred lines about 101\,Hz apart
that is present in every scan and unrelated to the event. And on 2018 April
15, fourteen months later, one OFF scan toward HIP\,99221 holds more than
twenty lines within 15\,kHz of the event's frequency, in a pattern that
repeats every 6\,kHz, all carrying their power in the same five time
samples: a pulsed, multi-line emitter. The event is different, a single
line on for minutes at a time on one steady drift, so this identifies a
human user of the frequency and not the event.

A third lead came from the review of this manuscript. The frequency lies in
2025--2110\,MHz, allocated to Earth-to-space and space-to-space links (and,
in the United States, shared with electronic news gathering). Spacecraft
transponders in this band pair an uplink with a downlink at 240/221 times
its frequency. The line's frequency times 240/221 is 2228.00007\,MHz at the
start of the second ON scan, and stays within 0.2\,kHz of 2228.000\,MHz
throughout the detection. A frequency chosen at random would fall that
close to a whole number of megahertz about once in 2500 times. This
suggests the uplink or forward-link carrier of a spacecraft whose downlink
is assigned 2228.000\,MHz, and such a carrier would also account for a line
that is on for minutes at a time and for a slow, steady drift, if the
transmitter compensates for Doppler shift. We found no published assignment
at 2228.0\,MHz to identify the spacecraft, and the hit tables hold nothing
within 50\,kHz of 2228.000\,MHz in that cadence. The coincidence was looked
for after the fact and is not direct evidence under our rules.

We consider a local or orbital transmitter the most likely origin, on the
allocation, the on-off behaviour and the transponder coincidence. The event
stays open. Several steps of the verification framework of
\citet{Sheikh2021} remain undone: asking the observatory's interference
group and consulting licence records for a transmitter near 2051.6\,MHz in
February 2017; testing the track for the curvature that the Earth's
rotation imposes on a sky source (about 1\,Hz over these 12 minutes,
comparable to our fit residual); and searching the scans for other lines
with the same on-off envelope and a drift proportional to frequency, since
our companion search required equal drift. The public products carry no
polarization information. The elevation was 76.5 degrees and the pointing
was 38 degrees from the geostationary belt.

\begin{figure*}
\plotone{figures/hip36357.png}
\caption{The six scans of the cadence of HIP\,36357 on 2017 February 23
around 2051.616734\,MHz: frequency across, time down, ON and OFF scans
alternating. Red triangles mark where the track of $-0.124$\,\hzs\ enters
and leaves each panel. The line is present in the first two ON scans only,
and switches off at the end of the first and on during the second.
\label{fig:hip36357}}
\end{figure*}

\section{Discussion} \label{sec:discussion}

\subsection{What the open list means}

The open list is where the method has least power, which is not the same as
where the signals are most interesting. Events stay open in crowded
spectra, where control levels are high and a line near any frequency is
common; at fast drifts, where the place of a line a few minutes later
cannot be predicted well and no test signal reaches; and when a line's
drift wanders between scans. Three of the five open flagged events have a
line of their drift toward other stars near their frequency and stay open
because that cannot be told from chance in those bands. Nineteen of the 51
open events of the queue do not exceed their control windows in any scan. A
reader interested in candidates should discount the open list accordingly;
a reader interested in the method should take it as a map of where forced
measurement against the strongest control fails.

The converse holds for closures. They are of unequal strength, and we have
tried to say which are weak: the 84 burst closures in which the line
survives the removal of its strongest sample, the five closures on
sidebands, the single-scan half-hour closures, and the two off-source
closures of flagged events that rest on a stacked margin just over the
threshold.

\subsection{What the re-examination confirms for searches}

None of the following is new as an idea; the re-examination measures how
much each is worth on this sample.

First, a forced measurement on the predicted track in every scan resolves
many cases that a threshold on independent detections leaves ambiguous: 19
of the 34 flagged events are measurably present off the source.
\citet{Price2020} note that signals below the detection threshold in the
OFF scans pass their filter, and visual inspection is step 2 of the
framework of \citet{Sheikh2021}. Second, time samples within a scan carry
information that a scan-integrated search discards
\citep[see also][]{Brzycki2023}: 450 of 771 near-miss hits, and 3 of the 34
flagged events, are dominated by a single 18\,s sample. Third, scans of
other targets hours or days away are a free control for local emitters, as
\citet{Price2020} and \citet{Painter2025} used them; a statistic over
frequency, and not over scans, turns a match into a probability and shows
where the test has no power.

The work also shows the limits of injection tests. Ours measured the rules
within their range and missed both of our serious errors, because the
injections shared the pipeline's assumptions, and they did not cover the
drifts of most of the sample. Known real signals, predictions checked on
data a rule had not seen, and an adversarial reading of the manuscript
against the code caught what the injections could not: the review found
that our first statistic for the test over days was invalid and that two
events we had called absent toward other stars were not.

\subsection{Limitations}

The sample is not random, and its selection makes two of the checks nearly
vacuous for it. Thresholds were chosen on small samples, and some on the
same controls that score them, so the quoted error rates are in-sample. No
test signal drifts faster than 0.5\,\hzs, and rules D5 to D9 were never run
on test signals. The burst rule closes on a 50\% share alone, which a
steady line under one unrelated burst would meet. The off-source test
missed two signals of \citet{Ma2023} that another pipeline finds off the
source. The statistic of Section~\ref{sec:probe} was designed after the
data were seen. The satellite comparison knows catalogued payloads only and
cannot see a satellite in a sidelobe. The work used only the L- and S-band
GBT products. Nothing here can confirm a signal: that needs new telescope
time.

\section{Conclusions} \label{sec:conclusions}

Using public data and partial file reads, we re-measured 771 near-miss hits
from the Breakthrough Listen 1327-star survey and the 34 events flagged for
follow-up in its inspection table, and found no candidate.

\begin{enumerate}
\item Of the 771 hits, 131 are shown to be local, 450 are dominated by a
single 18\,s sample and are not steady tones, 113 could not be reproduced,
and 70 hits in 51 events are not explained. Most of the unexplained events
lie where the method has little power.
\item Of the 34 flagged events, 29 are explained, 19 as measurable off the
source. Three of the 34 had been rejected by the survey's authors by
cross-reference to other stars, and the spectra confirm it.
\item A frequency-permutation statistic for the presence of an event's line
toward other stars over days explains three flagged events and makes
explicit that in crowded bands a line near the event's frequency is not
evidence.
\item The released hit files carry the drift-rate sign error noted in the
software's documentation; the survey's event table is corrected, the hit
files are not flagged.
\item The flagged events come mostly from low-declination targets observed
in mid-2016, a pattern these data cannot attribute.
\item One event, toward HIP\,36357 at 2051.617\,MHz, is unexplained. Its
frequency corresponds through the 240/221 transponder ratio to
2228.000\,MHz; a local or orbital transmitter is the most likely origin,
and site records could decide it.
\end{enumerate}

Within their range, the rules wrongly closed 0.4\% to 1.0\% of sky-only
test signals and caught 72\% of weak leaks. The code, the rules, the scripts
that produce every number and figure of this paper, the definitions of the
known-answer cases and the review reports are public at
\url{https://github.com/luizpersechini/cold-case-seti}; a table of the
outcome of every case accompanies the paper, and the per-case evidence is
available on request.

\section{Acknowledgments} \label{sec:ack}

This work rests entirely on the open data policy of Breakthrough Listen and
on the public products of \citet{Price2020} and \citet{Lebofsky2019}. It is
not affiliated with Breakthrough Listen, the SETI Institute or the
University of California, Berkeley.
%% AUTHOR: state here whether the findings were sent to the Breakthrough
%% Listen team, when, and what they replied.
The idea of donating idle computer time to this problem is that of
SETI@home \citep{Anderson2002,Anderson2025}.

The pipeline, its tests, the analysis scripts and a first draft of this
manuscript were written with a large-language-model coding assistant
(Claude Code, Anthropic, with Claude Opus and Fable models, November 2026).
Instances of Claude Sonnet served as the readers of Section~\ref{sec:ai},
and further instances reviewed the manuscript against the code and the
literature. Statements about prior work were checked against the cited
sources. The author directed the work, approved every closure made by hand,
and is responsible for the content.

\facility{GBT}

\software{turboSETI \citep{turboseti}, NumPy, Astropy, h5py, Matplotlib, DuckDB}

\bibliography{refs}{}
\bibliographystyle{aasjournal}

\end{document}
